\scanpage{93}
% archive foliation: 95
\uline{Courbes elliptiques réelles}
\deleted{Elles sont} La catégorie de ces courbes
(avec morphismes pas néc des iso) équivaut à celle
des systèmes $(\pi,\sigma,c)$, où $\pi$ est un $\mathbb Z$-module
libre, \inserted{de rg $2$}\deleted{\unsure{muni}} $c$ une structure \deleted{conf} conforme orientée
sur $\pi_{\mathbb R}=\pi\otimes_{\mathbb Z}\mathbb R$, et $\sigma$ un anti-homomorphisme du
syst. $(\pi,c)$, qui peut s'interpréter comme une
involution de $\pi$, renversant l'\uline{orientation}
\inserted{de $\pi\otimes\mathbb R$ i.e. de déterminant $-1$}
(donc $\ne\pm\operatorname{id}_\pi$),
et \unsure{conservant} la structure conforme. Pour le couple
$(\pi,\sigma)$ il y a exactement deux possibilités, à isom. près.
\deleted{Pour} (elles sont d'ailleurs isogènes\ldots). Il reste, pour chacune
d'elles, à trouver les structures conformes qui \unsure{est|sont}
invariantes par $\sigma$. On peut les normaliser par \deleted{$\unsure{\sigma\circ\sigma}$}
\[
Q(\unsure{e_+})\deleted{(=Q(\unsure{e_-}))}=1
\]
\margin{\uline{NB} $e_+$ est une base de \deleted{\unsure{k}}$\pi_+$\par
$e_-$\rule{3cm}{0.3pt}$\pi_-$}
donc ce sont les structures de la forme
\[
Q(x\unsure{e_+}+y\unsure{e_-})=x^2+ay^2
\]
où $a\in\mathbb R^{*+}$. Il y a donc deux familles de courbes
elliptiques réelles, paramétrées chacune par $\mathbb R^{*+}$ ; on peut
les désigner par $E_a,E'_a$ ($a\in\mathbb R^{*+}$). \deleted{Considérons}

Considérons les autos \deleted{de ces courbes}\inserted{d'une courbe elliptique}
définie par $(\pi,\sigma,c)$, ils correspondent aux \deleted{\unsure{Aut de}}
autos $\theta$ de $\pi$ qui a) commutent à $\sigma$ et b) respectent la
structure conforme $c$, \inserted{\uline{orientée}} \deleted{\unsure{en particulier}}
ce qui revient \unsure{au|à} $=$ \inserted{c'est à d' $\theta$ transforme $\pi^+_{\unsure{\pm}}$ en $\pi^+$}
la forme $Q$ (à un facteur près).
\inserted{(\unsure{en plus de C'est un $U$|en plus de l'isom.})}
\deleted{\unsure{Montrer}\unsure{ le}\unsure{ sg}\unsure{ fixe}\unsure{ est}\unsure{ le}}
Il transforme \deleted{\unsure{puisque}} $e_0$ en un multiple \inserted{ent} $\pm e_0$,
SPS quitte à composer avec la symétrie $\unsure{\lambda}=-\operatorname{id}$ que c'est $e_0$.
Mais le seul auto de $\mathbb R e_0+\mathbb R e_1$ qui respecte $Q$, qui fixe $e_0$ et qui
n'est pas l'identité, est $\sigma$, \deleted{\unsure{par conséquent}}\inserted{\unsure{exclu}} par l'orientation !
Donc il n'y a que l'\unsure{autom.} universel, \unsure{commun à tous} : $-\operatorname{id}_{\unsure{\pi}}$.
